\documentclass[11pt]{article}
\usepackage[margin=1in]{geometry}
\usepackage[T1]{fontenc}
\usepackage[utf8]{inputenc}
\usepackage{lmodern}
\usepackage{booktabs}
\usepackage{amsmath}
\usepackage{xcolor}
\usepackage[hidelinks]{hyperref}
\usepackage{natbib}
\setcitestyle{authoryear,round}

\newcommand{\pending}{\textcolor{red}{\textbf{[pending]}}}

\title{Where Was the Model Looking? Cheap Claim-Level Hallucination Detection\\ from Attention and Token Probabilities}
\author{Anonymous submission}
\date{}

\begin{document}
\maketitle

\begin{abstract}
Claim-level hallucination detection for long-form LLM output is far from solved: on the
MUCH benchmark the strongest published training-free method, Claim-Conditioned Probability
(CCP), reaches 0.77 ROC-AUC while spending more compute than generation itself. We study
GRASP, a detector that reads three interpretable attention statistics (grounding to the
prompt, attention-sink leakage and dispersion) from a few label-free-selected heads per layer,
pools them over relative depth, adds the token probabilities the model already produced, and
scores each claim with a small gradient-boosted classifier. On MUCH's human-annotated test
split GRASP reaches 0.924 ROC-AUC against CCP's 0.770 ($\Delta$ +0.154, 95\% CI
+0.110 to +0.199), wins on each of four open models, and costs 33--38\% of generation time
on a T4, where CCP measured on the same hardware costs 192--430\%. Retrained with the same
frozen recipe on RAGTruth it again beats CCP (0.827 vs 0.623). We report the limits as
plainly: the detector is supervised; trained on MUCH and applied unchanged to RAGTruth it only
ties CCP, even after label-free domain adaptation, and there the authors' pre-trained UHeads
transfer better (0.680 vs 0.639); its training-free variants do not beat CCP. With in-domain
labels GRASP also outperforms pre-trained UHeads on both benchmarks.
\end{abstract}

\section{Introduction}
An LLM states invented facts as fluently as true ones. For long answers, a single
answer-level confidence score does not tell a user \emph{which} sentence to check, so the useful
unit is the atomic claim \citep{min-etal-2023-factscore}. The MUCH benchmark
\citep{dentan2026muchmultilingualclaimhallucination} makes this setting measurable: 4,873
generations from four open models in four languages, pre-segmented into 20,751 claims, with the
generated tokens and their top-24 probabilities released. Its strongest baseline, CCP
\citep{fadeeva-etal-2024-fact}, filters token-level uncertainty through a natural-language
inference model applied at every token; it reaches 0.772 ROC-AUC at a cost the benchmark authors
put above that of generation.

We ask whether the model's own attention, read in one teacher-forced pass, can do better and
cheaper. The intuition is that when a model states a fact it holds, some heads retrieve from the
prompt; when it improvises, attention drifts to its own recent output or collapses onto the
attention sink. We make three contributions:
\begin{itemize}
  \item A lightweight claim-level detector, GRASP, built from three interpretable attention
        statistics on label-free-selected heads plus free token probabilities, and a small tree model.
  \item An evaluation that follows MUCH's protocol exactly, keeps every modelling choice on the
        training split, and reports per model, across models, on a second benchmark (RAGTruth;
        \citealp{niu-etal-2024-ragtruth}) and with runtime measured against CCP on the same GPU.
  \item An explicit account of what does \emph{not} work: zero-shot cross-benchmark transfer and
        training-free variants.
\end{itemize}

\section{Related work}
\textbf{Probability-based.} Token likelihood and entropy are cheap but conflate lexical and factual
uncertainty; CCP \citep{fadeeva-etal-2024-fact} removes lexical variation with NLI at a large
cost. \textbf{Sampling-based} methods such as SelfCheckGPT \citep{manakul-etal-2023-selfcheckgpt}
and semantic entropy \citep{kuhn2023semanticuncertaintylinguisticinvariances} multiply generation
cost and miss consistent errors. \textbf{Internal-state probes} predict hallucination from hidden
states \citep{azaria2023internalstatellmknows,chen2024insidellmsinternalstates} or attention:
Lookback Lens \citep{chuang2024lookbacklensdetectingmitigating} uses the ratio of attention to
context versus generated text, SinkProbe \citep{binkowski2026attentionsinksinternalsignals} the
attention-sink mass, and RAUQ \citep{vazhentsev2026efficienthallucinationdetectionllms} selects
uncertainty-aware heads without labels at sequence level. Closest to us, UHeads
\citep{shelmanov-etal-2025-head} train a transformer head on attention features and token
probabilities for claim-level detection. GRASP is in the same family; it differs in using three
interpretable statistics from a few label-free-selected heads and a small tree model, and in being
evaluated head-to-head against CCP on MUCH. Consistent with our choice of a small learner, \citet{lee2026meanshift} find that simple regularised probes on hidden states match far more complex ones.

\section{Method}
\subsection{Per-token attention signals}
The generated tokens are replayed through the model once. For generated token $t$ and a kept head,
with attention row $a_t$ over positions $0..t$ (position 0 is the sink, $1..P{-}1$ the prompt),
we compute grounding to the prompt $G_t$, sink leakage $S_t$ and dispersion $D_t$:
\begin{align}
G_t &= \frac{\sum_{j=1}^{P-1} a_{t,j}}{\sum_{j\ge 1} a_{t,j}}, \qquad S_t = a_{t,0},\\
D_t &= \frac{-\sum_{j\ge1} \tilde a_{t,j}\log \tilde a_{t,j}}{\log |\{j\ge1\}|},
\end{align}
where $\tilde a$ is the non-sink row renormalised. The row used is the one that \emph{produced}
token $t$. A saliency weight $W_t$ is the attention token $t$ receives from later generated tokens,
divided by their number.

\subsection{Label-free head selection}
For each layer we keep the $k$ heads whose $G$ varies most across tokens on 200 unlabelled
generations of the same model; a head with near-constant $G$ cannot discriminate. We use $k{=}4$.

\subsection{Claim features and classifier}
Layers are pooled into 8 relative-depth bins so that models of different depth share one feature
space. For each bin, kept head and signal, a claim's content tokens (stopwords and punctuation
removed) give the mean, max, min and saliency-weighted mean of the per-(layer, head) z-scored value
(384 features). We add seven token-probability features from the top-24 probabilities MUCH
already releases (chosen-token surprisal, entropy, top-1 probability, top-1/top-2 gap) and two
claim-length features. A gradient-boosted tree classifier (31 leaves, 800 rounds) predicts the
hallucination label. \textbf{GRASP-free}, a training-free variant, averages the per-model rank of
mid-depth minimum grounding with the rank of token surprisal and entropy, with signs fixed a priori.

\section{Experimental setup}
\textbf{Data.} MUCH, all 8 generation configs (Llama-3.2-3B, Llama-3.1-8B, Gemma-3-4B and
Ministral-8B, at temperatures 0.7 and 1.0). MUCH's test split is its 200 human-annotated generations
(867 claims); the train split is the other 4,673 (19,884 claims). We follow MUCH's evaluator:
GPT-4o labels on claims where GPT-4o and GPT-4.1 agree, the trailing end-of-sequence chunk dropped,
and stopwords masked for probability baselines.

\textbf{Protocol.} Feature groups, normalisation, learner and $k$ were chosen by 5-fold
cross-validation on the train split, grouped by generation and ranked by the \emph{mean of
per-model} ROC-AUCs, since pooled ROC can be inflated by recognising which model wrote a claim
(hallucination rates range from 22\% to 54\%). The selection was frozen to a file before the test
split was read. This was the second reading of MUCH's test split; the first evaluated an earlier
version (\S\ref{sec:limits}).

\textbf{Baselines.} CCP as released by MUCH ($\delta{=}24$, $\sigma{=}8$, \texttt{product}
aggregation), which we reproduce at 0.775 ROC / 0.643 PR (published: 0.772 / 0.639); token
surprisal and entropy; and UHeads \citep{shelmanov-etal-2025-head} via the authors' pre-trained
claim heads (\S\ref{sec:uheads}).

\textbf{Statistics.} 95\% CIs come from 2,000 bootstrap resamples of generations; differences
are paired on identical claims.

\section{Results}
\subsection{MUCH}
\begin{table}[h]
\centering\small
\begin{tabular}{lccccc}
\toprule
Method & Labels & ROC-AUC [95\% CI] & PR-AUC & TPR@10\% & $\Delta$ROC vs CCP \\
\midrule
CCP & no & 0.770 [0.724, 0.817] & 0.637 & 0.537 & -- \\
GRASP-0 (hand-weighted) & no & 0.629 [0.569, 0.687] & 0.388 & 0.188 & $-$0.141 \\
GRASP-free & no & 0.764 [0.713, 0.815] & 0.501 & 0.354 & $-$0.006 [$-$0.056, +0.041] \\
\textbf{GRASP} & yes & \textbf{0.924} [0.896, 0.948] & \textbf{0.845} & \textbf{0.790} & \textbf{+0.154} [+0.110, +0.199] \\
GRASP + CCP input & yes & 0.920 [0.895, 0.944] & 0.838 & 0.729 & +0.150 [+0.111, +0.191] \\
\bottomrule
\end{tabular}
\caption{MUCH human-annotated test split (867 claims, 229 hallucinated).}
\label{tab:main}
\end{table}

GRASP beats CCP on every model: Gemma-3-4B 0.917 vs 0.726, Llama-3.1-8B 0.937 vs 0.831,
Llama-3.2-3B 0.889 vs 0.767 and Ministral-8B 0.897 vs 0.830; all paired CIs exclude zero. Adding CCP
as an input does not help, so its NLI step is unnecessary. At response level (does an answer contain
any hallucination; 200 generations) GRASP reaches 0.911 vs CCP's 0.762 ($\Delta$ +0.149,
CI +0.086 to +0.214).

\textbf{Transfer to an unseen model.} Trained on three models' training claims and scored on the
fourth model's test claims, GRASP beats CCP on Gemma (0.808 vs 0.726), Llama-3.1-8B (0.882 vs 0.831)
and Ministral (0.853 vs 0.830) and trails it on Llama-3.2-3B (0.759 vs 0.767).

\textbf{What contributes.} In train cross-validation (mean per-model ROC; CCP 0.797), token
probabilities alone give 0.771, attention alone 0.832, both 0.866, and the tuned learner 0.885.
Four heads per layer beat one consistently (+0.01 to +0.02).

\subsection{Abstention}
Withholding the claims GRASP rates riskiest lowers the hallucination rate of what remains from
26.4\% to 11.4\% at 20\% withheld and to 7.2\% at 30\% (CCP: 16.3\% and 14.7\%); the area under
the risk--coverage curve is 0.062 vs 0.152.

\subsection{Second benchmark: RAGTruth}
We apply the frozen recipe to RAGTruth's Mistral-7B-Instruct responses (retrieval-augmented QA,
summarisation, data-to-text). Claims come from MUCH's segmenter and are hallucinated if they overlap
an annotated span; the test set has 445 responses and 17,226 claims (7.7\% hallucinated). Rules were
fixed before extraction (Table~\ref{tab:rag}).
\begin{table}[h]
\centering\small
\begin{tabular}{llccc}
\toprule
Method & Trained on & ROC-AUC [95\% CI] & PR-AUC & $\Delta$ROC vs CCP \\
\midrule
CCP & -- & 0.623 [0.602, 0.644] & 0.150 & -- \\
mean token entropy & -- & 0.623 & 0.114 & $-$0.000 \\
GRASP-free & -- & 0.595 & 0.098 & $-$0.028 [$-$0.057, +0.001] \\
GRASP, zero-shot & MUCH & 0.625 [0.606, 0.645] & 0.118 & +0.002 [$-$0.025, +0.029] \\
GRASP, zero-shot + rank norm. & MUCH + unlabelled RAGTruth & 0.639 [0.617, 0.659] & 0.131 & +0.015 [$-$0.011, +0.042] \\
\textbf{GRASP} & RAGTruth train & \textbf{0.827} [0.803, 0.849] & \textbf{0.383} & \textbf{+0.203} [+0.177, +0.233] \\
\bottomrule
\end{tabular}
\caption{RAGTruth test, claim level.}
\label{tab:rag}
\end{table}
Retrained in-domain, GRASP wins on every task (QA 0.859, summary 0.765, data-to-text 0.853, vs
CCP 0.626, 0.617, 0.624). Zero-shot from MUCH it only ties CCP: better on QA and summaries, worse
on data-to-text. As a label-free remedy we compared, by MUCH leave-one-model-out on the training
split, the frozen recipe (mean held-out ROC 0.808), per-domain rank normalisation of every feature
fitted on unlabelled target features (0.819) and rank normalisation with importance weighting
(0.818). The selected rank normalisation lifts zero-shot RAGTruth ROC to 0.639, still not
significantly above CCP.

\subsection{Comparison with UHeads}
\label{sec:uheads}
We ran the authors' pre-trained claim-level UHeads \citep{shelmanov-etal-2025-head}, teacher-forced
on exactly the claims GRASP is scored on. Our reimplementation of their attention features was
checked equal to their extractor. The score's sign and, on RAGTruth, the choice between their two
Mistral heads were fixed on training claims only.
\begin{table}[h]
\centering\small
\begin{tabular}{llcc}
\toprule
Benchmark & Method (labels) & ROC-AUC [95\% CI] & $\Delta$ROC vs UHead \\
\midrule
MUCH, Llama-3.1-8B & UHead (pre-trained) & 0.570 [0.486, 0.653] & -- \\
(282 claims) & CCP & 0.831 & +0.261 \\
 & GRASP (other three models) & 0.883 [0.817, 0.938] & +0.313 [+0.234, +0.397] \\
 & GRASP (in-domain) & \textbf{0.937} & +0.367 [+0.289, +0.449] \\
\midrule
RAGTruth, Mistral & CCP & 0.623 & $-$0.057 \\
(17,181 claims) & GRASP zero-shot + rank norm. (MUCH) & 0.639 & $-$0.041 [$-$0.076, $-$0.004] \\
 & UHead (pre-trained) & 0.680 [0.649, 0.708] & -- \\
 & GRASP (RAGTruth train) & \textbf{0.827} & +0.147 [+0.116, +0.180] \\
\bottomrule
\end{tabular}
\caption{GRASP vs pre-trained UHeads on identical claims.}
\label{tab:uheads}
\end{table}
With in-domain labels GRASP outperforms UHeads on both benchmarks. Zero-shot, UHeads' pre-trained
Mistral head transfers to RAGTruth better than GRASP trained on MUCH alone; it was trained on far
larger and more varied data. The pre-trained Llama head is weak on MUCH even on training claims
(oriented ROC 0.526), while the same code is sensible on RAGTruth. We therefore attribute this to a
mismatch between its training claims and MUCH's very short claims (about 4.6 tokens), and treat the
MUCH margin over UHeads as indicative only.

\subsection{Runtime}
On Tesla T4s, GRASP's extraction pass costs 33--38\% of measured generation time. MUCH's CCP,
timed on the same hardware, costs 235--430\% on MUCH's models and 192\% on RAGTruth-length
answers. CCP's NLI cost is fixed per token, so it is relatively worst for small, fast models.

\section{Limitations}
\label{sec:limits}
GRASP is supervised and needs labels from the setting it monitors; without them it transfers
across models but not across benchmarks and task families. Training-free variants do not beat
CCP at claim level on either benchmark, and at response level the RAGTruth advantage of the
training-free score (0.679 vs 0.647) does not hold on MUCH (0.616 vs 0.762). Without in-domain
labels, pre-trained UHeads transfer better than GRASP (RAGTruth 0.680 vs 0.639), and our
comparison uses their released heads rather than retraining them on MUCH. MUCH's test split was
read twice: once by an earlier version and once by GRASP as reported, with all choices frozen on
the training split beforehand. Test labels are GPT-4o's, applied to the human-annotated samples, as
in MUCH's evaluator. CCP and generation were timed on different sample generations, so overheads
compare per token. RAGTruth responses were replayed through Mistral-7B-Instruct-v0.1, assumed to
be the generator version. RAGTruth's original span labels miss many hallucinations \citep{yu2026rt4chart}, so absolute RAGTruth scores are likely depressed for every method.

\section{Conclusion}
Reading where the model looked, through a few interpretable attention statistics, plus the
probabilities it already produced, gives a claim-level hallucination detector that clearly beats
CCP and pre-trained UHeads on MUCH and, retrained, on RAGTruth, at a fraction of CCP's cost. The
open problem it leaves is transfer to new task families without labels, where large pre-trained
heads currently do better.

\bibliographystyle{plainnat}
\bibliography{../papers/references}
\end{document}
